\documentclass[11pt,a4paper]{article}
\usepackage{prof}
\ProfHeader{Two-state quantum model}

\begin{document}
\ProfTitle{Constructing a two-state Hamiltonian}
  {Build the matrix from the operator's action, then test whether a state has
   a definite energy.}

This short example demonstrates the document layout and linked practice.
For a new teaching document, replace its subject matter and choose the scope
from the learner's goal. It is a starter, not a complete quantum mechanics
course.

\section{What the matrix represents}

Assume that the orthonormal states $|A\rangle$ and $|B\rangle$ form the whole
space of our two-state model. The inner product measures the overlap of
two states; a state's norm squared is its inner product with itself.
``Orthonormal'' means that
$\langle A|A\rangle=\langle B|B\rangle=1$ and
$\langle A|B\rangle=0$. The first condition fixes each basis state's scale;
the second makes them independent, distinguishable basis alternatives.
In this ordered basis a state is
\[
 |\psi\rangle=c_A|A\rangle+c_B|B\rangle
 \quad\longleftrightarrow\quad
 \begin{pmatrix}c_A\\c_B\end{pmatrix}.
\]
The column records the state's coefficients. It is a representation of the
state in this chosen basis, not a new physical state.
For a normalized state, $|c_A|^2+|c_B|^2=1$: the squared coefficient
magnitudes are the probabilities of the two corresponding basis-state
measurement outcomes, which must add to one.

The Hamiltonian $\hat H$ is the linear energy operator. Linearity means
$\hat H(c_A|A\rangle+c_B|B\rangle)
=c_A\hat H|A\rangle+c_B\hat H|B\rangle$.
Therefore its action on the two basis states tells us its action on every
state in this model. Suppose that
\begin{align}
 \hat H|A\rangle&=(2\,\mathrm{eV})|A\rangle
                  -(1\,\mathrm{eV})|B\rangle,\\
 \hat H|B\rangle&=-(1\,\mathrm{eV})|A\rangle
                  +(2\,\mathrm{eV})|B\rangle.
\end{align}
Here eV is an energy unit, the electronvolt. Each coefficient in these
operator actions has units of energy.

\begin{profworked}{Read the output coefficients}
Multiplying a matrix by $(1,0)^{\mathsf T}$ selects its first column. Since
that input represents $|A\rangle$, the first column must be $(2,-1)^{\mathsf T}$
in eV. The second column must similarly represent $\hat H|B\rangle$.
Thus
\[
 H=\begin{pmatrix}2&-1\\-1&2\end{pmatrix}\,\mathrm{eV}.
\]
As a quick check, multiplying this matrix by $(1,0)^{\mathsf T}$ reproduces
the first stated operator action. An off-diagonal entry means that applying
the energy operator to one basis state produces a component along the other.
It is an energy-valued coupling, not a probability or a transition rate.
\end{profworked}

\section{Use the matrix to make a decision}

A definite-energy state satisfies $\hat H|\psi\rangle=E|\psi\rangle$:
the output is the original state multiplied by a single energy $E$.
This is the eigenvalue equation. The matrix representation gives the same
test using its coefficient column.

For the normalized state $(|A\rangle+|B\rangle)/\sqrt{2}$, direct multiplication
gives
\[
 \begin{pmatrix}2&-1\\-1&2\end{pmatrix}
 \frac{1}{\sqrt{2}}\begin{pmatrix}1\\1\end{pmatrix}
 =\frac{1}{\sqrt{2}}\begin{pmatrix}1\\1\end{pmatrix}.
\]
Restoring the matrix's eV unit gives $E=1\,\mathrm{eV}$. Both coefficients
are multiplied by the same number, so the state does have a definite energy.
The factor $1/\sqrt{2}$ normalizes the state because
$|1/\sqrt{2}|^2+|1/\sqrt{2}|^2=1$.

\begin{profattempt}{changed-site}{Does the same state still work?}
The local environment of state $|B\rangle$ changes. The new operator actions are
\[
 \hat H|A\rangle=(2|A\rangle-|B\rangle)\,\mathrm{eV},\qquad
 \hat H|B\rangle=(-|A\rangle+4|B\rangle)\,\mathrm{eV}.
\]
Construct the new matrix and use it to decide whether
$(|A\rangle+|B\rangle)/\sqrt{2}$ still has a definite energy.
Give one calculation that justifies your decision. This is one connected
task: represent the changed model and interpret its action.
\end{profattempt}

\begin{table}[htbp]
\centering
\begin{tabularx}{\linewidth}{@{}p{0.29\linewidth}Y@{}}
\toprule
If you are stuck on & Return to this idea \\
\midrule
Choosing matrix entries & Each column is the output for its corresponding
  input basis state. \\
Interpreting the result & A definite energy requires one common multiplier
  for the entire coefficient column. \\
\bottomrule
\end{tabularx}
\caption{Use the specific missing idea before repeating the whole explanation.}
\end{table}

The two-site equality used in the worked example was a property of that
model, not a rule for every two-state Hamiltonian. Changing the environment
tests whether the method still works when that property is removed.

\begin{profattempt}{relative-sign}{A later check}
After a gap, repeat the definite-energy decision with
\[
 H=\begin{pmatrix}3&-2\\-2&3\end{pmatrix}\,\mathrm{eV},\qquad
 |\psi\rangle=\frac{|A\rangle-|B\rangle}{\sqrt{2}}.
\]
Give the energy if it exists and justify your decision. The relative sign
between the two coefficients has changed; use the same operator-action
test rather than guessing from the earlier result.
\end{profattempt}

% Group hints after the lesson route, then put full solutions on a separate
% page so a hint link does not expose the complete answer beside the hint.
\clearpage
\section{Hints}

Use a hint if you cannot identify a next step, then return to the attempt.
The full solutions are on a separate page: follow the matching link when
you are ready to check the complete reasoning.

\begin{profhint}{changed-site}{Compare both coefficients}
The changed second basis-state action supplies the second column. After
multiplication by $(1,1)^{\mathsf T}/\sqrt{2}$, ask whether the first and
second output coefficients have the same ratio to their input coefficients.
\end{profhint}

\begin{profhint}{relative-sign}{Track the minus sign}
The input column is $(1,-1)^{\mathsf T}/\sqrt{2}$. Multiply both rows by
this column, then compare the whole output with the whole input.
\end{profhint}

\clearpage
\section{Full solutions}

After checking the reasoning, return to the attempt and reproduce its key
decision without reading the answer.

\begin{profsolution}{changed-site}{The symmetry assumption no longer applies}
The first and second columns give
\[
 H=\begin{pmatrix}2&-1\\-1&4\end{pmatrix}\,\mathrm{eV}.
\]
Applying it to the proposed state gives
\[
 H\frac{1}{\sqrt{2}}\begin{pmatrix}1\\1\end{pmatrix}
 =\frac{1}{\sqrt{2}}\begin{pmatrix}1\\3\end{pmatrix}\,\mathrm{eV}.
\]
The first coefficient would require $E=1\,\mathrm{eV}$, while the second
would require $E=3\,\mathrm{eV}$. There is no common multiplier. The proposed
state is therefore not a definite-energy state of this changed Hamiltonian.
Computing an average energy would answer a different question: an average
does not establish that the state has only one possible measured energy.

If your matrix was right but you accepted the state anyway, the point to
repair is the common-multiplier test. If your matrix was wrong, reconstruct
its columns from the stated basis-state actions before continuing.
\end{profsolution}

\begin{profsolution}{relative-sign}{One multiplier for both entries}
Direct multiplication gives
\[
 H\frac{1}{\sqrt{2}}\begin{pmatrix}1\\-1\end{pmatrix}
 =\frac{1}{\sqrt{2}}\begin{pmatrix}5\\-5\end{pmatrix}\,\mathrm{eV}
 =(5\,\mathrm{eV})\frac{1}{\sqrt{2}}\begin{pmatrix}1\\-1\end{pmatrix}.
\]
Both input coefficients are multiplied by $5\,\mathrm{eV}$, so the state
has definite energy $5\,\mathrm{eV}$. Opposite output signs are consistent
with one multiplier because the input signs were already opposite.
\end{profsolution}

\end{document}
